\section*{1. CONTEXT AND REQUIREMENTS}

\phantomsection
\addcontentsline{toc}{section}{1. CONTEXT AND REQUIREMENTS}

\subsection*{1.1. TMC Context and Problem Statement}

\phantomsection
\addcontentsline{toc}{subsection}{1.1. TMC Context and Problem Statement}

TMC Manufacturing Joint Stock Company operates in the manufacturing sector and has two factories in different provinces: Binh Duong with about 800 employees and Long An with about 600 employees. TMC wants one system that controls who can enter the workplace and, at the same time, collects the data needed for monthly attendance reports.

Every employee identifies themselves with a QR code shown on a personal smartphone. At each lane, a kiosk reads the code and the system checks whether this employee may enter the factory. The gate itself is not new: TMC already has a turnstile system, and our system has to open and close the gate through its API. Each entry or exit then becomes an attendance transaction.

The attendance data does not stay inside the system. At the end of each month it is sent to TMC's existing HR/Payroll system through an API, and it has to be available for reconciliation before 12:00 PM on the first day of the next month.

\Needspace{9\baselineskip}
\subsection*{1.2. System Objectives}

\phantomsection
\addcontentsline{toc}{subsection}{1.2. System Objectives}

From this context, we set four objectives for the system. The first three follow the main business needs, and the last one groups the quality targets given in the assignment.

\begin{itemize}[leftmargin=36pt, itemsep=2pt, topsep=4pt, parsep=2pt]
  \item \textbf{OBJ-01 -- Access control:} Identify employees by QR code and let only authorized employees pass the existing turnstiles.
  \item \textbf{OBJ-02 -- Attendance recording:} Record each entry and exit as an attendance transaction, without loss or duplication.
  \item \textbf{OBJ-03 -- Reporting and handover:} Provide monthly attendance reports and transfer the data to HR/Payroll before noon on the first day of each month.
  \item \textbf{OBJ-04 -- Quality:} Meet the assignment's requirements for kiosk response time, no record loss, RBAC with audit logging, and personal-data protection.
\end{itemize}

\Needspace{9\baselineskip}
\subsection*{1.3. System Scope}

\phantomsection
\addcontentsline{toc}{subsection}{1.3. System Scope}

\subsubsection*{1.3.1. In Scope}

\phantomsection
\addcontentsline{toc}{subsubsection}{1.3.1. In Scope}

The current phase covers both the gate side and the office side of attendance:

\begin{itemize}[leftmargin=36pt, itemsep=2pt, topsep=4pt, parsep=2pt]
  \item Employee information and QR codes.
  \item Stations, lanes, and kiosks at both factories.
  \item Gate control through the existing turnstile system.
  \item Entry and exit recording, with review and adjustment of attendance data.
  \item Attendance queries and monthly reports.
  \item Accounts, roles, permissions, and login.
  \item Audit logs, operational monitoring, and administration.
  \item Delivery of attendance data to HR/Payroll.
\end{itemize}

\Needspace{9\baselineskip}
\subsubsection*{1.3.2. Out of Scope}

\phantomsection
\addcontentsline{toc}{subsubsection}{1.3.2. Out of Scope}

Two areas are excluded by the assignment for this phase:

\begin{itemize}[leftmargin=36pt, itemsep=2pt, topsep=4pt, parsep=2pt]
  \item \textbf{Visitor management:} Visitors and their entries and exits are not handled.
  \item \textbf{Payroll calculation:} Salaries, allowances, and deductions stay in HR/Payroll. Our system only provides the attendance data.
\end{itemize}

\Needspace{9\baselineskip}
\subsection*{1.4. Stakeholders, Actors and External Systems}

\phantomsection
\addcontentsline{toc}{subsection}{1.4. Stakeholders, Actors and External Systems}

We first describe the people who use the system directly, then the other parties who care about its results, and finally the systems and devices it has to work with.

\subsubsection*{1.4.1. User Groups}

\phantomsection
\addcontentsline{toc}{subsubsection}{1.4.1. User Groups}

The assignment lists six user groups. Their main tasks in our system are:

\begin{itemize}[leftmargin=36pt, itemsep=2pt, topsep=4pt, parsep=2pt]
  \item \textbf{Employee:} Shows a fresh QR code on a smartphone, passes through the gates, views their own attendance records, and reports discrepancies.
  \item \textbf{Production Supervisor:} Monitors attendance in the assigned area and helps verify exceptional cases, but cannot modify attendance directly.
  \item \textbf{HR Staff:} Maintains employee data, access permissions, and shift schedules. HR staff also review discrepancies, approve adjustments with evidence, and monitor the HR/Payroll handover.
  \item \textbf{Security Guard:} Monitors gate results, assists employees who cannot show a QR code, and opens a gate in an emergency within their authority.
  \item \textbf{System Administrator:} Manages accounts, roles, and station, lane, and kiosk settings, and monitors, backs up, and restores the system. Business data is not accessible to this role by default.
  \item \textbf{Executive Management:} Views reports and aggregated statistics.
\end{itemize}

\Needspace{9\baselineskip}
\subsubsection*{1.4.2. Other Stakeholders}

\phantomsection
\addcontentsline{toc}{subsubsection}{1.4.2. Other Stakeholders}

Some parties never log in, but the design still has to respect what they need:

\begin{itemize}[leftmargin=36pt, itemsep=2pt, topsep=4pt, parsep=2pt]
  \item \textbf{TMC (system owner):} Needs reliable attendance data, an on-time monthly handover, and an acceptable operating cost.
  \item \textbf{Employees as data subjects:} Their personal data should be used only for access control and attendance.
  \item \textbf{Operators of the turnstile and HR/Payroll systems:} Need clear API contracts and agreed responsibilities when a failure occurs.
\end{itemize}

\Needspace{9\baselineskip}
\subsubsection*{1.4.3. External Systems and Employee Channels}

\phantomsection
\addcontentsline{toc}{subsubsection}{1.4.3. External Systems and Employee Channels}

Two of the systems around ours already exist at TMC and cannot be changed by us. The others are devices and channels that employees and staff use every day.

\begin{itemize}[leftmargin=36pt, itemsep=2pt, topsep=4pt, parsep=2pt]
  \item \textbf{Turnstile control system:} An existing system. The assignment states that it offers APIs to open and close the gate for each access event, but not which signals it reports back, so we state them as assumptions I-01 and I-05.
  \item \textbf{HR/Payroll system:} TMC's existing payroll system. It receives attendance data from our system through its API (IC-01).
  \item \textbf{Kiosk:} One kiosk per lane, with a QR scanner and a display. The kiosk software is part of our system.
  \item \textbf{Smartphone, mobile app, web portal, email:} The channels named in the assignment for displaying a QR code. How email is used is explained in assumption I-08.
  \item \textbf{Management interface:} Used by guards, HR staff, supervisors, administrators, and executives, each according to their role.
\end{itemize}

\Needspace{9\baselineskip}
\subsection*{1.5. Constraints from the Assignment}

\phantomsection
\addcontentsline{toc}{subsection}{1.5. Constraints from the Assignment}

A few points are fixed by the assignment and are not design choices: employees are identified by QR code, every employee has a personal smartphone to display it, and TMC currently has two factories, Binh Duong and Long An. The assignment also requires two integrations. It does not say what message formats, response signals, or availability these APIs have, so our assumptions about them are collected in Table 1.1.

\begin{itemize}[leftmargin=36pt, itemsep=2pt, topsep=4pt, parsep=2pt]
  \item \textbf{IC-01 -- HR/Payroll integration:} Attendance data is sent to the existing HR/Payroll system through its API and must be available for reconciliation before 12:00 PM on the first day of each month. Payroll calculation stays in HR/Payroll.
  \item \textbf{IC-02 -- Turnstile integration:} Gates are opened and closed through the existing turnstile system's API. Because gate opening is part of the kiosk response, this call counts toward the 2-second limit of NFR-01.
\end{itemize}

The M2 assignment adds two issues that were not in M1: loss of connectivity during a transaction, and exposure of static QR codes. We answer them in Section 5.8, but we mention them here because their solutions affect the architecture from the start.

\subsection*{1.6. Functional Requirements}

\phantomsection
\addcontentsline{toc}{subsection}{1.6. Functional Requirements}

The assignment only names five functional areas, so the numbered sub-requirements below are our own breakdown of each area. Most of them restate what the area name already implies. Five of them go beyond the assignment text and come from our analysis of the context; they are marked \textit{(added)}: FR-01.2, FR-02.4, FR-04.4, FR-05.2, and FR-05.4.

\subsubsection*{1.6.1. FR-01 -- User and QR Code Management}

\phantomsection
\addcontentsline{toc}{subsubsection}{1.6.1. FR-01 -- User and QR Code Management}

\begin{itemize}[leftmargin=36pt, itemsep=2pt, topsep=4pt, parsep=2pt]
  \item \textbf{FR-01.1:} The system shall manage the information needed to identify employees, control access, and record attendance.
  \item \textbf{FR-01.2 \textit{(added)}:} The system shall issue each employee a short-lived QR code whose authenticity can be checked and which carries no unnecessary personal information.
  \item \textbf{FR-01.3:} The system shall let employees obtain and display QR codes on their smartphones through the mobile app or web portal, and obtain a new code when the current one expires. Company email may link to the authenticated QR screen (I-08).
  \item \textbf{FR-01.4:} The system shall let authorized users activate, suspend, revoke, and reinstate an employee's access permission.
\end{itemize}

\subsubsection*{1.6.2. FR-02 -- Gate Station Management}

\phantomsection
\addcontentsline{toc}{subsubsection}{1.6.2. FR-02 -- Gate Station Management}

\begin{itemize}[leftmargin=36pt, itemsep=2pt, topsep=4pt, parsep=2pt]
  \item \textbf{FR-02.1:} The system shall manage stations, lanes, and kiosks at each factory, including whether a lane is for entry or exit.
  \item \textbf{FR-02.2:} The system shall open and close the gate through the existing turnstile system for each authorized entry or exit. It shall keep apart four outcomes: a received command, a confirmed opening, a confirmed passage, and a confirmed closing.
  \item \textbf{FR-02.3:} The system shall monitor kiosks and controllers and report faults that affect entry and exit.
  \item \textbf{FR-02.4 \textit{(added)}:} The system shall let authorized users open a gate in an emergency and record every such opening.
\end{itemize}

\subsubsection*{1.6.3. FR-03 -- Attendance Event Recording}

\phantomsection
\addcontentsline{toc}{subsubsection}{1.6.3. FR-03 -- Attendance Event Recording}

\begin{itemize}[leftmargin=36pt, itemsep=2pt, topsep=4pt, parsep=2pt]
  \item \textbf{FR-03.1:} The system shall receive QR codes scanned at kiosks and return the result to the kiosk.
  \item \textbf{FR-03.2:} For a QR scan, the system shall allow a gate to open only when the QR code is valid, the employee is active, and the employee has access permission for that factory. Guard-assisted access (I-04) and emergency opening (FR-02.4) follow their own authorization rules.
  \item \textbf{FR-03.3:} The system shall record an attendance transaction for each confirmed entry or exit. Records created through guard assistance require supporting evidence and approval.
  \item \textbf{FR-03.4:} The system shall not create a duplicate transaction when the same scan is resubmitted.
  \item \textbf{FR-03.5:} Each transaction shall include the employee ID, the timestamp in Vietnam time, the direction, the factory, the lane, and the source of the record.
  \item \textbf{FR-03.6:} The system shall retain rejected, failed, and undetermined scans for traceability.
\end{itemize}

\subsubsection*{1.6.4. FR-04 -- Reporting and Auditing}

\phantomsection
\addcontentsline{toc}{subsubsection}{1.6.4. FR-04 -- Reporting and Auditing}

\begin{itemize}[leftmargin=36pt, itemsep=2pt, topsep=4pt, parsep=2pt]
  \item \textbf{FR-04.1:} The system shall provide access and attendance monitoring information to authorized users.
  \item \textbf{FR-04.2:} The system shall produce monthly attendance reports for each employee, factory, and period.
  \item \textbf{FR-04.3:} The system shall transfer attendance data to HR/Payroll through its API in the format HR/Payroll requires.
  \item \textbf{FR-04.4 \textit{(added)}:} The system shall support reviewing and adjusting attendance through a request and approval process that preserves the original data.
  \item \textbf{FR-04.5:} The system shall maintain audit logs of administrative actions, permission changes, gate results, attendance adjustments, and handovers.
  \item \textbf{FR-04.6:} The system shall allow employees to view their own attendance records, and supervisors and HR staff to view records within their assigned scope.
  \item \textbf{FR-04.7:} The system shall let authorized users search and view audit logs.
\end{itemize}

\subsubsection*{1.6.5. FR-05 -- System Administration}

\phantomsection
\addcontentsline{toc}{subsubsection}{1.6.5. FR-05 -- System Administration}

\begin{itemize}[leftmargin=36pt, itemsep=2pt, topsep=4pt, parsep=2pt]
  \item \textbf{FR-05.1:} The system shall manage accounts, roles, and permissions.
  \item \textbf{FR-05.2 \textit{(added)}:} The system shall let HR staff manage shift schedules and attendance aggregation rules.
  \item \textbf{FR-05.3:} The system shall monitor operational status and raise alerts for faults in kiosks, controllers, or the HR/Payroll integration.
  \item \textbf{FR-05.4 \textit{(added)}:} The system shall support backup and recovery of its data.
  \item \textbf{FR-05.5:} The system shall authenticate users, reject invalid logins, allow logout, and restrict functions according to the user's role.
\end{itemize}

\Needspace{9\baselineskip}
\subsection*{1.7. Non-functional Requirements}

\phantomsection
\addcontentsline{toc}{subsection}{1.7. Non-functional Requirements}

The assignment gives four non-functional requirements. We keep their targets unchanged, but each one is short, so after the target we add how we read it in our design. For NFR-01 this is where the response time is measured; for NFR-03 and NFR-04 it is what ``comprehensive'' logging and data protection mean for this system.

\begin{itemize}[leftmargin=36pt, itemsep=2pt, topsep=4pt, parsep=2pt]
  \item \textbf{NFR-01 -- Performance:} Kiosk response time shall not exceed 2 seconds at shift change, with about 800 entry and exit events in 30 minutes per factory. We measure it from when the kiosk reads the QR code until it displays the result.
  \item \textbf{NFR-02 -- Reliability:} Attendance records shall not be lost during normal system operation.
  \item \textbf{NFR-03 -- Security and auditability:} The system shall enforce RBAC with complete audit logging. In our reading, every administrative action and every entry or exit attempt, including rejected ones, is logged.
  \item \textbf{NFR-04 -- Privacy:} Personal data of employees shall be protected as required by the applicable data protection law. In our reading, access is limited by role, and each component stores only the personal data it needs for its responsibilities.
\end{itemize}

\Needspace{9\baselineskip}
\subsubsection*{1.7.1. Shift-change Workload Estimate}

\phantomsection
\addcontentsline{toc}{subsubsection}{1.7.1. Shift-change Workload Estimate}

To get a feel for the load behind NFR-01, we turned the workload in the assignment into rates. At a shift change, each factory has about 800 events within 30 minutes (1,800 seconds), which gives an average of
\[
\frac{800}{1800} \approx 0.44 \mbox{ events/second} \;(\approx 26.7 \mbox{ events/minute})
\]
per factory. When both factories change shifts at the same time, the shared components see
\[
\frac{2 \times 800}{1800} \approx 0.89 \mbox{ events/second}.
\]

These numbers are averages, and real arrivals are not spread evenly over 30 minutes. We expect most employees to arrive in the last few minutes before the shift starts, so the short peaks are higher. This is our assumption, not a figure from the assignment; Section 2.1.3 turns it into a burst rate, and NFR-01 still has to hold for every scan during these peaks.

\Needspace{9\baselineskip}
\subsection*{1.8. Assumptions and Team Policy}

\phantomsection
\addcontentsline{toc}{subsection}{1.8. Assumptions and Team Policy}

While reading the assignment, we found several questions it does not answer. What does the turnstile API report back? Does HR/Payroll confirm that it received the data? How long should a QR code stay valid, and how can email show one? Our design depends on these answers, so Table 1.1 records the answer we use for each question, and later chapters refer to it by ID. The Status column shows whether a point still has to be checked with the API owner or approved by TMC.

{\small\renewcommand{\arraystretch}{1.08}
\begin{longtable}{|>{\raggedright\arraybackslash}p{28pt}|>{\raggedright\arraybackslash}p{325pt}|>{\raggedright\arraybackslash}p{64pt}|}
\hline
\multicolumn{1}{|>{\columncolor{tableheader}\centering\arraybackslash}p{28pt}|}{\textbf{ID}} & \multicolumn{1}{>{\columncolor{tableheader}\centering\arraybackslash}p{325pt}|}{\textbf{Assumption or policy}} & \multicolumn{1}{>{\columncolor{tableheader}\centering\arraybackslash}p{64pt}|}{\textbf{Status}} \\
\hline
\endhead
\multicolumn{3}{c}{\rule{0pt}{20pt}\normalsize\textbf{Table 1.1. Assumptions, Team Policy, and Confirmation Status}} \\
\endlastfoot
I-01 & \textbf{Controller contract.} The controller API distinguishes request receipt from confirmed opening (OPEN\_CONFIRMED) and reports a passage event with a stable identifier linked to the lane and command. & Verify API \\
\hline
I-02 & \textbf{HR/Payroll receipt.} HR/Payroll receives data in batches and returns a receipt acknowledgement or a clear rejection of the content. A handover succeeds only when this acknowledgement is received; without either answer the outcome is undetermined. & Verify API \\
\hline
I-03 & \textbf{No blind resend.} An attempt with an undetermined controller outcome is not treated as successful, and the opening command is not resent automatically. & Team policy \\
\hline
I-04 & \textbf{Assisted access.} A security guard assists an employee who cannot present a QR code, or whose QR scan cannot be completed, after verifying identity and current access eligibility. The attempt counts toward attendance only after HR approval. & TMC to approve \\
\hline
I-05 & \textbf{Closing contract.} The controller reports LANE\_CLEAR only when the opening opportunity has ended and the lane is physically clear. The system then sends closeGate once, and the controller's safety interlock reports CLOSED. If no passage arrives within 15 seconds, the case goes to reconciliation; this timer never authorizes closing. & Verify API \\
\hline
I-06 & \textbf{Shift-independent access.} Factory access does not depend on the work-shift schedule; shifts govern attendance aggregation only. & Open \\
\hline
I-07 & \textbf{Passage-based attendance.} A normal QR attendance transaction is created only after the I-01 passage event is durably recorded. Opening alone never creates attendance; missing passage evidence goes to reconciliation. & Team policy \\
\hline
I-08 & \textbf{Email channel.} Company email leads the employee to an authenticated, fresh QR display instead of carrying a reusable QR image. & Open \\
\hline
I-09 & \textbf{QR policy.} A signed QR code is valid for 60 seconds, is bound to one factory, and can be claimed by only one ACCEPTED scan across that factory's lanes. & Team policy \\*
\hline\noalign{\nobreak}
I-10 & \textbf{Eligibility freshness and degraded mode.} Factory-local access-right data counts as fresh while its last source-confirmed checkpoint is at most 60 seconds old. After that, the factory may enter degraded mode and keep accepting valid signed QR codes for up to 30 minutes after that checkpoint, flagging each transaction for review. It stops earlier if another required condition fails, such as the clock-error bound or a durable write (Section 5.7.2). TMC has to accept that an employee revoked during this window may still pass; the review flag cannot undo that entry. & TMC to approve \\
\hline
\end{longtable}
}
\phantomsection\addcontentsline{lot}{table}{\protect\numberline{1.1}Assumptions, Team Policy, and Confirmation Status}


\Needspace{9\baselineskip}
\subsection*{1.9. Main Business Flows}

\phantomsection
\addcontentsline{toc}{subsection}{1.9. Main Business Flows}

Two business flows matter most for the architecture: the daily one at the gate and the monthly one towards HR/Payroll. We describe them here only at the business level; the detailed runtime views are in Section 4.2.

\subsubsection*{1.9.1. Gate Entry/Exit and Attendance Recording Flow}

\phantomsection
\addcontentsline{toc}{subsubsection}{1.9.1. Gate Entry/Exit and Attendance Recording Flow}

In the normal case, an employee passes the gate as follows:

\begin{enumerate}[leftmargin=36pt, itemsep=2pt, topsep=4pt, parsep=2pt]
  \item The employee opens a fresh QR code on the smartphone and presents it at the kiosk.
  \item The kiosk reads the QR code, and the system validates it and checks access permission.
  \item If access is allowed and the lane is free, the system reserves the lane, claims the QR code, and stores the accepted attempt in a single write. Only after this write does it send the opening command.
  \item The controller confirms opening. The system stores this result, and the kiosk displays it within 2 seconds.
  \item When the controller reports passage, the system records the attendance transaction.
  \item When the controller reports LANE\_CLEAR, the system closes the gate.
\end{enumerate}

Steps 4 to 6 do not wait for each other. After the confirmed opening is stored, the kiosk can show its result, while attendance and closing are handled whenever the controller sends the passage and LANE\_CLEAR signals.

\Needspace{6\baselineskip}
The normal case is only part of the story. The flow also has to handle these exceptions:

\begin{itemize}[leftmargin=36pt, itemsep=2pt, topsep=4pt, parsep=2pt]
  \item an invalid QR code, or an employee without permission;
  \item a busy lane;
  \item access data that is too old;
  \item a resubmitted scan;
  \item a controller error or no response from the controller;
  \item a missing passage or a missing closing confirmation;
  \item an employee who has no QR code.
\end{itemize}

\Needspace{9\baselineskip}
\subsubsection*{1.9.2. Attendance Data Handover to HR/Payroll Flow}

\phantomsection
\addcontentsline{toc}{subsubsection}{1.9.2. Attendance Data Handover to HR/Payroll Flow}

The handover runs once a month, after the month has closed:

\begin{enumerate}[leftmargin=36pt, itemsep=2pt, topsep=4pt, parsep=2pt]
  \item Aggregate attendance data for the previous month.
  \item Resolve missing records, discrepancies, and undetermined outcomes within the applicable authority.
  \item Package the data in the format HR/Payroll requires and send it through the API.
  \item Receive the HR/Payroll acknowledgement.
  \item If the handover is not confirmed, follow the receipt and resubmission rules of Flow C (Section 4.2.4). Later adjustments are sent as supplementary handovers.
\end{enumerate}
